\chapter{Software design views and function data flow}\label{app:design-views}
This appendix translates the parser-generated function index into software design views. The delivered application is a single offline HTML file, so ``client'', ``controller/service'' and ``repository'' are responsibility boundaries inside one script, not separate deployed processes. The optional Java boundary is shown only as an integration point for a future host; no Java call is made by the current page.

\begin{infobox}
The diagrams use UML-style sequence and activity conventions for readable design views. Each function in Appendix~\ref{app:index} receives a flow record in Table~\ref{tab:function-flow-map}: its exact HTML lines, responsibility boundary, inputs, outputs and next flow. This keeps the schematic complete without drawing one box per symbol.
\end{infobox}

\section{Composition and data-flow view}
\begin{figure}[htbp]\centering
\begin{tikzpicture}[font=\footnotesize,box/.style={draw,rounded corners,align=center,minimum height=12mm,text width=3.3cm,inner sep=4pt,fill=#1!8,draw=#1!60!black},box/.default=blue,arr/.style={-{Stealth[length=2mm]},thick,gray!70!black}]
\node[box=gray] (client) at (0,0) {Browser client / DOM, buttons, tabs, SVG/canvas};
\node[box=blue] (svc) at (5,0) {Controller / service-style functions / decode, SOP, charts, review};
\node[box=orange] (repo) at (10,0) {Repository / data access / File bytes, ZIP, hashes, exports};
\node[box=green] (model) at (5,-2.7) {RUNS[] + history + profile / common in-memory model};
\node[box=red] (events) at (10,-2.7) {Findings and status / structured events, notices, limits};
\draw[arr] (client)--node[above]{commands / selections}(svc);
\draw[arr] (svc)--node[above]{read / write}(repo);
\draw[arr] (repo)--node[right]{decoded values}(model);
\draw[arr] (model)--(svc);
\draw[arr] (svc)--node[above]{render data}(client);
\draw[arr] (svc)--(events); \draw[arr] (events)--(client);
\end{tikzpicture}
\caption{Responsibility view of the current single-file page.}\label{fig:design-composition}
\end{figure}

\section{Intake and analysis activity}
\begin{figure}[htbp]\centering
\begin{tikzpicture}[font=\footnotesize,box/.style={draw,rounded corners,align=center,minimum height=10mm,text width=3.0cm,inner sep=3pt,fill=#1!8,draw=#1!60!black},box/.default=blue,arr/.style={-{Stealth[length=2mm]},thick,gray!70!black}]
\node[box=gray](f) at (0,0){.ixo / .eds / pasted table};
\node[box=blue](i) at (4,0){intake() / readStaged()};
\node[box=orange](d) at (8,0){decodeIxo() / decodeEds()};
\node[box=green](m) at (12,0){common RUNS[] / protocol, wells, curves};
\node[box=violet](s) at (4,-2.4){sopEvaluateRun()/procedure + interpretation};
\node[box=red](q) at (8,-2.4){reviewForensicEvents() / rising curves, integrity, findings};
\node[box=orange](c) at (12,-2.4){ccCompute() / control limits and signals};
\draw[arr](f)--(i)--(d)--(m);\draw[arr](m)--(s);\draw[arr](m)--(q);\draw[arr](m)--(c); \draw[arr](s)--(q);\draw[arr](q)--(c);
\end{tikzpicture}
\caption{Activity view from staged files to method interpretation, forensic review and control charts.}\label{fig:design-activity}
\end{figure}

\section{Client-to-service sequence}
\begin{figure}[htbp]\centering
\begin{tikzpicture}[font=\scriptsize,box/.style={draw,rounded corners,align=center,minimum height=8mm,text width=2.5cm,inner sep=3pt,fill=#1!8,draw=#1!60!black},box/.default=blue,arr/.style={-{Stealth[length=2mm]},thick,gray!70!black}]
\node[box=gray](u) at (0,0){Operator / button};\node[box=blue](v) at (3.6,0){DOM + view state};\node[box=orange](x) at (7.2,0){service function};\node[box=green](r) at (10.8,0){RUNS / history};\node[box=red](e) at (14.4,0){finding / export};
\draw[arr](u)--node[above]{select / click}(v);\draw[arr](v)--node[above]{call}(x);\draw[arr](x)--node[above]{read / derive}(r);\draw[arr](r)--node[above]{rows}(x);\draw[arr](x)--node[above]{render / download}(e);\draw[arr](e)--node[below]{status}(v);
\end{tikzpicture}
\caption{Synchronous local sequence for a button action. No network or Java endpoint participates in the current implementation.}\label{fig:design-sequence}
\end{figure}

\section{Optional host integration boundary}
If a future Java host embeds the page, the boundary is an adapter around file selection and export delivery. The host may provide a file stream and receive a generated artifact, while the decoder, SOP engine, chart engine and evidence model remain deterministic page services. The contract should preserve the common run model and the profile hash; it must not rewrite stored Cq values or calls.

\section{Function-level flow records}
Table~\ref{tab:function-flow-map} is generated from the same current HTML snapshot as Appendix~\ref{app:index}; a changed function name or line range therefore appears in both places on the next build.
\begin{landscape}\scriptsize
\begin{longtable}{@{}p{3.2cm}p{1.45cm}p{2.7cm}p{4.0cm}p{4.0cm}p{3.0cm}@{}}
\caption{Function-level data-flow records for the current HTML snapshot.}\label{tab:function-flow-map}\\
\toprule Function & Lines & Responsibility & Inputs & Outputs & Next flow\\\midrule\endfirsthead
\toprule Function & Lines & Responsibility & Inputs & Outputs & Next flow\\\midrule\endhead
ZH\_TERMS & 996--1920 & Shared model / utility & primitive values or shared state & normalised value or configuration & utility → caller \\\hline
zhPhrase & 1924--2002 & Shared model / utility & primitive values or shared state & normalised value or configuration & utility → caller \\\hline
ui & 2003--2003 & Client / view & DOM state, selection, user command & HTML/SVG, status and download command & view state → render \\\hline
localizeDOM & 2004--2023 & Client / view & DOM state, selection, user command & HTML/SVG, status and download command & view state → render \\\hline
showCuteToast & 2024--2030 & Client / view & DOM state, selection, user command & HTML/SVG, status and download command & view state → render \\\hline
languageStorageKey & 2031--2031 & Client / view & DOM state, selection, user command & HTML/SVG, status and download command & view state → render \\\hline
rememberedLanguage & 2032--2032 & Client / view & DOM state, selection, user command & HTML/SVG, status and download command & view state → render \\\hline
setLanguage & 2033--2043 & Client / view & DOM state, selection, user command & HTML/SVG, status and download command & view state → render \\\hline
ROLES & 2045--2045 & Shared model / utility & primitive values or shared state & normalised value or configuration & utility → caller \\\hline
REP\_REGEX & 2046--2050 & Shared model / utility & primitive values or shared state & normalised value or configuration & utility → caller \\\hline
CALL & 2052--2052 & Shared model / utility & primitive values or shared state & normalised value or configuration & utility → caller \\\hline
SAMPLE\_TYPE & 2053--2057 & Shared model / utility & primitive values or shared state & normalised value or configuration & utility → caller \\\hline
TYPE\_TO\_ROLE & 2058--2062 & Shared model / utility & primitive values or shared state & normalised value or configuration & utility → caller \\\hline
NEGATIVE\_FAMILY & 2063--2063 & Shared model / utility & primitive values or shared state & normalised value or configuration & utility → caller \\\hline
ROLE\_COLOURS & 2064--2065 & Shared model / utility & primitive values or shared state & normalised value or configuration & utility → caller \\\hline
ROLE\_LINE & 2066--2067 & Shared model / utility & primitive values or shared state & normalised value or configuration & utility → caller \\\hline
PALETTE & 2068--2068 & Shared model / utility & primitive values or shared state & normalised value or configuration & utility → caller \\\hline
ce & 2070--2070 & Shared model / utility & primitive values or shared state & normalised value or configuration & utility → caller \\\hline
clone & 2077--2077 & Shared model / utility & primitive values or shared state & normalised value or configuration & utility → caller \\\hline
esc & 2078--2078 & Shared model / utility & primitive values or shared state & normalised value or configuration & utility → caller \\\hline
xesc & 2079--2079 & Shared model / utility & primitive values or shared state & normalised value or configuration & utility → caller \\\hline
safeName & 2080--2080 & Shared model / utility & primitive values or shared state & normalised value or configuration & utility → caller \\\hline
bytesFrom & 2081--2081 & Repository / data access & File bytes, decoded rows or export request & stored model, checksum or file bytes & repository → model/export \\\hline
scrollToEl & 2082--2082 & Client / view & DOM state, selection, user command & HTML/SVG, status and download command & view state → render \\\hline
option & 2083--2083 & Client / view & DOM state, selection, user command & HTML/SVG, status and download command & view state → render \\\hline
num & 2084--2084 & Shared model / utility & primitive values or shared state & normalised value or configuration & utility → caller \\\hline
safeRe & 2085--2085 & Shared model / utility & primitive values or shared state & normalised value or configuration & utility → caller \\\hline
uniq & 2086--2086 & Shared model / utility & primitive values or shared state & normalised value or configuration & utility → caller \\\hline
byKey & 2087--2087 & Shared model / utility & primitive values or shared state & normalised value or configuration & utility → caller \\\hline
mean & 2092--2092 & Shared model / utility & primitive values or shared state & normalised value or configuration & utility → caller \\\hline
sd & 2094--2094 & Shared model / utility & primitive values or shared state & normalised value or configuration & utility → caller \\\hline
median & 2095--2095 & Shared model / utility & primitive values or shared state & normalised value or configuration & utility → caller \\\hline
percentile & 2097--2102 & Shared model / utility & primitive values or shared state & normalised value or configuration & utility → caller \\\hline
range & 2103--2103 & Shared model / utility & primitive values or shared state & normalised value or configuration & utility → caller \\\hline
cv & 2104--2104 & Shared model / utility & primitive values or shared state & normalised value or configuration & utility → caller \\\hline
lgamma & 2107--2112 & Shared model / utility & primitive values or shared state & normalised value or configuration & utility → caller \\\hline
betacf & 2114--2133 & Shared model / utility & primitive values or shared state & normalised value or configuration & utility → caller \\\hline
betai & 2134--2138 & Shared model / utility & primitive values or shared state & normalised value or configuration & utility → caller \\\hline
tDistTwoTail & 2140--2143 & Shared model / utility & primitive values or shared state & normalised value or configuration & utility → caller \\\hline
gammap & 2145--2160 & Shared model / utility & primitive values or shared state & normalised value or configuration & utility → caller \\\hline
chiSqP & 2161--2161 & Shared model / utility & primitive values or shared state & normalised value or configuration & utility → caller \\\hline
tTest & 2164--2187 & Shared model / utility & primitive values or shared state & normalised value or configuration & utility → caller \\\hline
fitLin & 2188--2195 & Shared model / utility & primitive values or shared state & normalised value or configuration & utility → caller \\\hline
CRC\_TABLE & 2200--2200 & Client / view & DOM state, selection, user command & HTML/SVG, status and download command & view state → render \\\hline
crc32 & 2201--2201 & Repository / data access & File bytes, decoded rows or export request & stored model, checksum or file bytes & repository → model/export \\\hline
sourceSha256 & 2206--2212 & Repository / data access & File bytes, decoded rows or export request & stored model, checksum or file bytes & repository → model/export \\\hline
sha256Hex & 2213--2245 & Repository / data access & File bytes, decoded rows or export request & stored model, checksum or file bytes & repository → model/export \\\hline
inflate & 2246--2251 & Repository / data access & File bytes, decoded rows or export request & stored model, checksum or file bytes & repository → model/export \\\hline
readZip & 2253--2278 & Repository / data access & File bytes, decoded rows or export request & stored model, checksum or file bytes & repository → model/export \\\hline
dosStamp & 2279--2283 & Shared model / utility & primitive values or shared state & normalised value or configuration & utility → caller \\\hline
makeStoredZip & 2284--2309 & Repository / data access & File bytes, decoded rows or export request & stored model, checksum or file bytes & repository → model/export \\\hline
propText & 2317--2317 & Repository / data access & File bytes, decoded rows or export request & stored model, checksum or file bytes & repository → model/export \\\hline
pv & 2318--2318 & Shared model / utility & primitive values or shared state & normalised value or configuration & utility → caller \\\hline
propValues & 2319--2319 & Repository / data access & File bytes, decoded rows or export request & stored model, checksum or file bytes & repository → model/export \\\hline
base64ToBytes & 2320--2320 & Repository / data access & File bytes, decoded rows or export request & stored model, checksum or file bytes & repository → model/export \\\hline
posToWell & 2321--2321 & Shared model / utility & primitive values or shared state & normalised value or configuration & utility → caller \\\hline
posToLimsWell & 2322--2322 & Shared model / utility & primitive values or shared state & normalised value or configuration & utility → caller \\\hline
wellToPos & 2323--2326 & Shared model / utility & primitive values or shared state & normalised value or configuration & utility → caller \\\hline
hexBEToDouble & 2327--2334 & Shared model / utility & primitive values or shared state & normalised value or configuration & utility → caller \\\hline
decodeAcq & 2336--2423 & Controller / service & RUNS, profile, history or chart inputs & derived rows, findings, limits or calls & model → next service/view \\\hline
acquisitionStats & 2424--2431 & Client / view & DOM state, selection, user command & HTML/SVG, status and download command & view state → render \\\hline
decodePlateModel & 2434--2495 & Client / view & DOM state, selection, user command & HTML/SVG, status and download command & view state → render \\\hline
decodeSubsets & 2497--2507 & Controller / service & RUNS, profile, history or chart inputs & derived rows, findings, limits or calls & model → next service/view \\\hline
decodeProtocol & 2509--2544 & Controller / service & RUNS, profile, history or chart inputs & derived rows, findings, limits or calls & model → next service/view \\\hline
ANALYSIS\_KINDS & 2552--2560 & Controller / service & RUNS, profile, history or chart inputs & derived rows, findings, limits or calls & model → next service/view \\\hline
analysisKind & 2561--2565 & Controller / service & RUNS, profile, history or chart inputs & derived rows, findings, limits or calls & model → next service/view \\\hline
LC\_CALC\_STATE\_LABELS & 2570--2570 & Controller / service & RUNS, profile, history or chart inputs & derived rows, findings, limits or calls & model → next service/view \\\hline
calcStateLabel & 2571--2571 & Controller / service & RUNS, profile, history or chart inputs & derived rows, findings, limits or calls & model → next service/view \\\hline
cpMethodOf & 2572--2578 & Shared model / utility & primitive values or shared state & normalised value or configuration & utility → caller \\\hline
RESULT\_PROPS & 2580--2581 & Repository / data access & File bytes, decoded rows or export request & stored model, checksum or file bytes & repository → model/export \\\hline
readLists & 2582--2590 & Repository / data access & File bytes, decoded rows or export request & stored model, checksum or file bytes & repository → model/export \\\hline
resultKindOf & 2591--2601 & Shared model / utility & primitive values or shared state & normalised value or configuration & utility → caller \\\hline
PAIRING\_RULES & 2608--2608 & Shared model / utility & primitive values or shared state & normalised value or configuration & utility → caller \\\hline
rqResult & 2609--2625 & Shared model / utility & primitive values or shared state & normalised value or configuration & utility → caller \\\hline
decodeRelQuant & 2626--2672 & Controller / service & RUNS, profile, history or chart inputs & derived rows, findings, limits or calls & model → next service/view \\\hline
ownedBy & 2677--2685 & Shared model / utility & primitive values or shared state & normalised value or configuration & utility → caller \\\hline
ownQueryAll & 2686--2686 & Shared model / utility & primitive values or shared state & normalised value or configuration & utility → caller \\\hline
ownQuery & 2687--2687 & Shared model / utility & primitive values or shared state & normalised value or configuration & utility → caller \\\hline
decodeAnalyses & 2689--2776 & Controller / service & RUNS, profile, history or chart inputs & derived rows, findings, limits or calls & model → next service/view \\\hline
ownsResult & 2783--2791 & Shared model / utility & primitive values or shared state & normalised value or configuration & utility → caller \\\hline
analysisResults & 2792--2811 & Controller / service & RUNS, profile, history or chart inputs & derived rows, findings, limits or calls & model → next service/view \\\hline
decodeIxo & 2813--3155 & Controller / service & RUNS, profile, history or chart inputs & derived rows, findings, limits or calls & model → next service/view \\\hline
EDS & 3160--3708 & Shared model / utility & primitive values or shared state & normalised value or configuration & utility → caller \\\hline
\$ & 3163--3163 & Shared model / utility & primitive values or shared state & normalised value or configuration & utility → caller \\\hline
esc & 3164--3164 & Shared model / utility & primitive values or shared state & normalised value or configuration & utility → caller \\\hline
num & 3165--3165 & Shared model / utility & primitive values or shared state & normalised value or configuration & utility → caller \\\hline
fmt & 3166--3166 & Shared model / utility & primitive values or shared state & normalised value or configuration & utility → caller \\\hline
pill & 3167--3167 & Shared model / utility & primitive values or shared state & normalised value or configuration & utility → caller \\\hline
dt & 3168--3168 & Shared model / utility & primitive values or shared state & normalised value or configuration & utility → caller \\\hline
hex & 3169--3169 & Shared model / utility & primitive values or shared state & normalised value or configuration & utility → caller \\\hline
mean & 3170--3170 & Shared model / utility & primitive values or shared state & normalised value or configuration & utility → caller \\\hline
sd & 3171--3171 & Shared model / utility & primitive values or shared state & normalised value or configuration & utility → caller \\\hline
wellPos & 3173--3173 & Shared model / utility & primitive values or shared state & normalised value or configuration & utility → caller \\\hline
download & 3174--3178 & Repository / data access & File bytes, decoded rows or export request & stored model, checksum or file bytes & repository → model/export \\\hline
safeName & 3179--3179 & Shared model / utility & primitive values or shared state & normalised value or configuration & utility → caller \\\hline
CRC\_T & 3182--3182 & Repository / data access & File bytes, decoded rows or export request & stored model, checksum or file bytes & repository → model/export \\\hline
crc32 & 3183--3183 & Repository / data access & File bytes, decoded rows or export request & stored model, checksum or file bytes & repository → model/export \\\hline
MD5 & 3186--3213 & Shared model / utility & primitive values or shared state & normalised value or configuration & utility → caller \\\hline
inflateRaw & 3216--3221 & Repository / data access & File bytes, decoded rows or export request & stored model, checksum or file bytes & repository → model/export \\\hline
readZip & 3222--3249 & Repository / data access & File bytes, decoded rows or export request & stored model, checksum or file bytes & repository → model/export \\\hline
makeZip & 3252--3268 & Repository / data access & File bytes, decoded rows or export request & stored model, checksum or file bytes & repository → model/export \\\hline
xml & 3271--3271 & Shared model / utility & primitive values or shared state & normalised value or configuration & utility → caller \\\hline
kid & 3272--3272 & Shared model / utility & primitive values or shared state & normalised value or configuration & utility → caller \\\hline
kids & 3273--3273 & Shared model / utility & primitive values or shared state & normalised value or configuration & utility → caller \\\hline
txt & 3274--3274 & Shared model / utility & primitive values or shared state & normalised value or configuration & utility → caller \\\hline
argb & 3275--3275 & Shared model / utility & primitive values or shared state & normalised value or configuration & utility → caller \\\hline
bracketList & 3276--3276 & Shared model / utility & primitive values or shared state & normalised value or configuration & utility → caller \\\hline
ini & 3277--3279 & Shared model / utility & primitive values or shared state & normalised value or configuration & utility → caller \\\hline
FLAGS & 3284--3302 & Shared model / utility & primitive values or shared state & normalised value or configuration & utility → caller \\\hline
flagName & 3303--3303 & Shared model / utility & primitive values or shared state & normalised value or configuration & utility → caller \\\hline
parseEds & 3307--3513 & Controller / service & RUNS, profile, history or chart inputs & derived rows, findings, limits or calls & model → next service/view \\\hline
recalcCt & 3517--3521 & Controller / service & RUNS, profile, history or chart inputs & derived rows, findings, limits or calls & model → next service/view \\\hline
parseLog & 3523--3539 & Controller / service & RUNS, profile, history or chart inputs & derived rows, findings, limits or calls & model → next service/view \\\hline
studyRows & 3541--3545 & Shared model / utility & primitive values or shared state & normalised value or configuration & utility → caller \\\hline
QS\_DEFAULT\_COLS & 3553--3559 & Shared model / utility & primitive values or shared state & normalised value or configuration & utility → caller \\\hline
qsCols & 3560--3565 & Shared model / utility & primitive values or shared state & normalised value or configuration & utility → caller \\\hline
pad2 & 3566--3566 & Shared model / utility & primitive values or shared state & normalised value or configuration & utility → caller \\\hline
qsCalDate & 3567--3567 & Shared model / utility & primitive values or shared state & normalised value or configuration & utility → caller \\\hline
qsStamp & 3568--3572 & Shared model / utility & primitive values or shared state & normalised value or configuration & utility → caller \\\hline
qsRGB & 3573--3573 & Shared model / utility & primitive values or shared state & normalised value or configuration & utility → caller \\\hline
qsInstrument & 3576--3587 & Shared model / utility & primitive values or shared state & normalised value or configuration & utility → caller \\\hline
javaHashOrder & 3589--3593 & Shared model / utility & primitive values or shared state & normalised value or configuration & utility → caller \\\hline
qsHeader & 3594--3613 & Shared model / utility & primitive values or shared state & normalised value or configuration & utility → caller \\\hline
qsTables & 3614--3692 & Client / view & DOM state, selection, user command & HTML/SVG, status and download command & view state → render \\\hline
qsExport & 3693--3705 & Repository / data access & File bytes, decoded rows or export request & stored model, checksum or file bytes & repository → model/export \\\hline
QS\_DYE\_FILTER & 3721--3723 & Shared model / utility & primitive values or shared state & normalised value or configuration & utility → caller \\\hline
QS\_TASK\_ROLE & 3724--3725 & Shared model / utility & primitive values or shared state & normalised value or configuration & utility → caller \\\hline
QS\_AMP & 3726--3727 & Shared model / utility & primitive values or shared state & normalised value or configuration & utility → caller \\\hline
QS\_KIND & 3728--3729 & Shared model / utility & primitive values or shared state & normalised value or configuration & utility → caller \\\hline
isoLocal & 3730--3735 & Shared model / utility & primitive values or shared state & normalised value or configuration & utility → caller \\\hline
isEdsName & 3736--3736 & Shared model / utility & primitive values or shared state & normalised value or configuration & utility → caller \\\hline
decodeEds & 3737--3865 & Controller / service & RUNS, profile, history or chart inputs & derived rows, findings, limits or calls & model → next service/view \\\hline
ixoIntegrity & 3868--3882 & Repository / data access & File bytes, decoded rows or export request & stored model, checksum or file bytes & repository → model/export \\\hline
integrityLabel & 3883--3887 & Repository / data access & File bytes, decoded rows or export request & stored model, checksum or file bytes & repository → model/export \\\hline
integrityEvents & 3888--3903 & Controller / service & RUNS, profile, history or chart inputs & derived rows, findings, limits or calls & model → next service/view \\\hline
edsEvents & 3904--3987 & Controller / service & RUNS, profile, history or chart inputs & derived rows, findings, limits or calls & model → next service/view \\\hline
ec & 3988--3988 & Shared model / utility & primitive values or shared state & normalised value or configuration & utility → caller \\\hline
edsRecordTables & 3991--4036 & Client / view & DOM state, selection, user command & HTML/SVG, status and download command & view state → render \\\hline
drawTemperatureTrace & 4037--4049 & Client / view & DOM state, selection, user command & HTML/SVG, status and download command & view state → render \\\hline
edsDerived & 4052--4102 & Controller / service & RUNS, profile, history or chart inputs & derived rows, findings, limits or calls & model → next service/view \\\hline
pseudonymMap & 4106--4115 & Repository / data access & File bytes, decoded rows or export request & stored model, checksum or file bytes & repository → model/export \\\hline
pseudoOn & 4116--4116 & Shared model / utility & primitive values or shared state & normalised value or configuration & utility → caller \\\hline
exportName & 4117--4117 & Repository / data access & File bytes, decoded rows or export request & stored model, checksum or file bytes & repository → model/export \\\hline
pseudoRun & 4118--4125 & Controller / service & RUNS, profile, history or chart inputs & derived rows, findings, limits or calls & model → next service/view \\\hline
withExportNames & 4126--4130 & Repository / data access & File bytes, decoded rows or export request & stored model, checksum or file bytes & repository → model/export \\\hline
rqRows & 4131--4138 & Shared model / utility & primitive values or shared state & normalised value or configuration & utility → caller \\\hline
edsSignalRows & 4139--4155 & Shared model / utility & primitive values or shared state & normalised value or configuration & utility → caller \\\hline
xlsxCol & 4158--4158 & Repository / data access & File bytes, decoded rows or export request & stored model, checksum or file bytes & repository → model/export \\\hline
xlsxCell & 4159--4165 & Repository / data access & File bytes, decoded rows or export request & stored model, checksum or file bytes & repository → model/export \\\hline
makeXlsx & 4166--4180 & Repository / data access & File bytes, decoded rows or export request & stored model, checksum or file bytes & repository → model/export \\\hline
qsWorkbook & 4181--4188 & Shared model / utility & primitive values or shared state & normalised value or configuration & utility → caller \\\hline
reviewForensicEvents & 4190--4197 & Controller / service & RUNS, profile, history or chart inputs & derived rows, findings, limits or calls & model → next service/view \\\hline
SOP\_OUTCOMES & 4208--4208 & Controller / service & RUNS, profile, history or chart inputs & derived rows, findings, limits or calls & model → next service/view \\\hline
SOP\_OUTCOME\_DEFAULTS & 4209--4220 & Controller / service & RUNS, profile, history or chart inputs & derived rows, findings, limits or calls & model → next service/view \\\hline
sopBase & 4221--4237 & Controller / service & RUNS, profile, history or chart inputs & derived rows, findings, limits or calls & model → next service/view \\\hline
SOP\_PRESETS & 4238--4265 & Controller / service & RUNS, profile, history or chart inputs & derived rows, findings, limits or calls & model → next service/view \\\hline
STORED\_SOP & 4335--4335 & Controller / service & RUNS, profile, history or chart inputs & derived rows, findings, limits or calls & model → next service/view \\\hline
sopLoadStored & 4340--4340 & Controller / service & RUNS, profile, history or chart inputs & derived rows, findings, limits or calls & model → next service/view \\\hline
sopStore & 4341--4341 & Controller / service & RUNS, profile, history or chart inputs & derived rows, findings, limits or calls & model → next service/view \\\hline
sopMaybeSelectScDefault & 4342--4349 & Controller / service & RUNS, profile, history or chart inputs & derived rows, findings, limits or calls & model → next service/view \\\hline
sopNormalise & 4350--4364 & Controller / service & RUNS, profile, history or chart inputs & derived rows, findings, limits or calls & model → next service/view \\\hline
sopCanonical & 4365--4368 & Controller / service & RUNS, profile, history or chart inputs & derived rows, findings, limits or calls & model → next service/view \\\hline
sopHash & 4369--4369 & Controller / service & RUNS, profile, history or chart inputs & derived rows, findings, limits or calls & model → next service/view \\\hline
sopJSON & 4370--4370 & Controller / service & RUNS, profile, history or chart inputs & derived rows, findings, limits or calls & model → next service/view \\\hline
sopChanged & 4371--4371 & Controller / service & RUNS, profile, history or chart inputs & derived rows, findings, limits or calls & model → next service/view \\\hline
sopMatcher & 4374--4380 & Controller / service & RUNS, profile, history or chart inputs & derived rows, findings, limits or calls & model → next service/view \\\hline
sopNum & 4381--4381 & Controller / service & RUNS, profile, history or chart inputs & derived rows, findings, limits or calls & model → next service/view \\\hline
sopCq & 4382--4382 & Controller / service & RUNS, profile, history or chart inputs & derived rows, findings, limits or calls & model → next service/view \\\hline
sopOmitted & 4383--4383 & Controller / service & RUNS, profile, history or chart inputs & derived rows, findings, limits or calls & model → next service/view \\\hline
sopTargetSpec & 4384--4390 & Controller / service & RUNS, profile, history or chart inputs & derived rows, findings, limits or calls & model → next service/view \\\hline
sopRunApplicability & 4391--4401 & Controller / service & RUNS, profile, history or chart inputs & derived rows, findings, limits or calls & model → next service/view \\\hline
sopReplicateMinimum & 4402--4407 & Controller / service & RUNS, profile, history or chart inputs & derived rows, findings, limits or calls & model → next service/view \\\hline
CC\_TARGET\_ALIAS & 4410--4410 & Controller / service & RUNS, profile, history or chart inputs & derived rows, findings, limits or calls & model → next service/view \\\hline
ccCanonTarget & 4411--4411 & Controller / service & RUNS, profile, history or chart inputs & derived rows, findings, limits or calls & model → next service/view \\\hline
sopControlIdentity & 4412--4434 & Controller / service & RUNS, profile, history or chart inputs & derived rows, findings, limits or calls & model → next service/view \\\hline
sopLabScreening & 4435--4453 & Controller / service & RUNS, profile, history or chart inputs & derived rows, findings, limits or calls & model → next service/view \\\hline
ccBaselineNeed & 4454--4454 & Controller / service & RUNS, profile, history or chart inputs & derived rows, findings, limits or calls & model → next service/view \\\hline
ccNoDataReason & 4455--4460 & Controller / service & RUNS, profile, history or chart inputs & derived rows, findings, limits or calls & model → next service/view \\\hline
controlBatchRecords & 4461--4470 & Controller / service & RUNS, profile, history or chart inputs & derived rows, findings, limits or calls & model → next service/view \\\hline
controlBatchGraph & 4471--4487 & Client / view & DOM state, selection, user command & HTML/SVG, status and download command & view state → render \\\hline
sopControlSpec & 4489--4494 & Controller / service & RUNS, profile, history or chart inputs & derived rows, findings, limits or calls & model → next service/view \\\hline
sopFmt & 4495--4495 & Controller / service & RUNS, profile, history or chart inputs & derived rows, findings, limits or calls & model → next service/view \\\hline
sopTime & 4496--4500 & Controller / service & RUNS, profile, history or chart inputs & derived rows, findings, limits or calls & model → next service/view \\\hline
sopFmtTime & 4501--4502 & Controller / service & RUNS, profile, history or chart inputs & derived rows, findings, limits or calls & model → next service/view \\\hline
wellSignal & 4505--4510 & Shared model / utility & primitive values or shared state & normalised value or configuration & utility → caller \\\hline
runTimeline & 4513--4533 & Controller / service & RUNS, profile, history or chart inputs & derived rows, findings, limits or calls & model → next service/view \\\hline
runTimes & 4534--4541 & Controller / service & RUNS, profile, history or chart inputs & derived rows, findings, limits or calls & model → next service/view \\\hline
sopStdCurves & 4544--4557 & Controller / service & RUNS, profile, history or chart inputs & derived rows, findings, limits or calls & model → next service/view \\\hline
sopEvaluateRun & 4560--4696 & Controller / service & RUNS, profile, history or chart inputs & derived rows, findings, limits or calls & model → next service/view \\\hline
sopEvaluateAll & 4697--4701 & Controller / service & RUNS, profile, history or chart inputs & derived rows, findings, limits or calls & model → next service/view \\\hline
sopResultForWell & 4702--4705 & Controller / service & RUNS, profile, history or chart inputs & derived rows, findings, limits or calls & model → next service/view \\\hline
sopMeta & 4706--4706 & Controller / service & RUNS, profile, history or chart inputs & derived rows, findings, limits or calls & model → next service/view \\\hline
sopInterpretationRows & 4707--4718 & Controller / service & RUNS, profile, history or chart inputs & derived rows, findings, limits or calls & model → next service/view \\\hline
sopSampleRows & 4719--4732 & Controller / service & RUNS, profile, history or chart inputs & derived rows, findings, limits or calls & model → next service/view \\\hline
sopCriteriaRows & 4733--4736 & Controller / service & RUNS, profile, history or chart inputs & derived rows, findings, limits or calls & model → next service/view \\\hline
sopRunRows & 4737--4748 & Controller / service & RUNS, profile, history or chart inputs & derived rows, findings, limits or calls & model → next service/view \\\hline
sopEvents & 4749--4761 & Controller / service & RUNS, profile, history or chart inputs & derived rows, findings, limits or calls & model → next service/view \\\hline
invalidateAnalysisCaches & 4767--4769 & Controller / service & RUNS, profile, history or chart inputs & derived rows, findings, limits or calls & model → next service/view \\\hline
forensicSnapshot & 4770--4775 & Controller / service & RUNS, profile, history or chart inputs & derived rows, findings, limits or calls & model → next service/view \\\hline
extraForensicEvents & 4776--4776 & Controller / service & RUNS, profile, history or chart inputs & derived rows, findings, limits or calls & model → next service/view \\\hline
reviewForensicEventsCore & 4777--4777 & Controller / service & RUNS, profile, history or chart inputs & derived rows, findings, limits or calls & model → next service/view \\\hline
forensicCounts & 4778--4784 & Controller / service & RUNS, profile, history or chart inputs & derived rows, findings, limits or calls & model → next service/view \\\hline
extraForensicEventsUncached & 4785--4827 & Controller / service & RUNS, profile, history or chart inputs & derived rows, findings, limits or calls & model → next service/view \\\hline
sopOutcomeFor & 4832--4836 & Controller / service & RUNS, profile, history or chart inputs & derived rows, findings, limits or calls & model → next service/view \\\hline
SOP\_OUTCOME\_OPTIONS & 4837--4837 & Client / view & DOM state, selection, user command & HTML/SVG, status and download command & view state → render \\\hline
sopField & 4838--4847 & Controller / service & RUNS, profile, history or chart inputs & derived rows, findings, limits or calls & model → next service/view \\\hline
SOP\_ACTION\_LABEL & 4848--4850 & Controller / service & RUNS, profile, history or chart inputs & derived rows, findings, limits or calls & model → next service/view \\\hline
sopSetPath & 4851--4856 & Controller / service & RUNS, profile, history or chart inputs & derived rows, findings, limits or calls & model → next service/view \\\hline
sopTableEditor & 4857--4875 & Client / view & DOM state, selection, user command & HTML/SVG, status and download command & view state → render \\\hline
sopProfileEdited & 4876--4879 & Controller / service & RUNS, profile, history or chart inputs & derived rows, findings, limits or calls & model → next service/view \\\hline
renderSopHeader & 4880--4888 & Client / view & DOM state, selection, user command & HTML/SVG, status and download command & view state → render \\\hline
sopProcedureHTML & 4889--4908 & Controller / service & RUNS, profile, history or chart inputs & derived rows, findings, limits or calls & model → next service/view \\\hline
renderSopEditor & 4909--4999 & Client / view & DOM state, selection, user command & HTML/SVG, status and download command & view state → render \\\hline
SOP\_VIEW & 5000--5000 & Controller / service & RUNS, profile, history or chart inputs & derived rows, findings, limits or calls & model → next service/view \\\hline
renderSopResults & 5001--5048 & Client / view & DOM state, selection, user command & HTML/SVG, status and download command & view state → render \\\hline
sopCurrentRows & 5049--5052 & Client / view & DOM state, selection, user command & HTML/SVG, status and download command & view state → render \\\hline
csvOf & 5053--5053 & Repository / data access & File bytes, decoded rows or export request & stored model, checksum or file bytes & repository → model/export \\\hline
sheetOf & 5054--5054 & Shared model / utility & primitive values or shared state & normalised value or configuration & utility → caller \\\hline
sopReadme & 5055--5068 & Controller / service & RUNS, profile, history or chart inputs & derived rows, findings, limits or calls & model → next service/view \\\hline
sopDownloadWorkbook & 5069--5076 & Controller / service & RUNS, profile, history or chart inputs & derived rows, findings, limits or calls & model → next service/view \\\hline
sopDownloadBundle & 5077--5086 & Controller / service & RUNS, profile, history or chart inputs & derived rows, findings, limits or calls & model → next service/view \\\hline
wireSop & 5087--5103 & Client / view & DOM state, selection, user command & HTML/SVG, status and download command & view state → render \\\hline
renderSop & 5104--5104 & Client / view & DOM state, selection, user command & HTML/SVG, status and download command & view state → render \\\hline
renderResults & 5105--5105 & Client / view & DOM state, selection, user command & HTML/SVG, status and download command & view state → render \\\hline
svgEsc & 5108--5108 & Client / view & DOM state, selection, user command & HTML/SVG, status and download command & view state → render \\\hline
niceTicks & 5109--5116 & Shared model / utility & primitive values or shared state & normalised value or configuration & utility → caller \\\hline
fmtTick & 5117--5117 & Shared model / utility & primitive values or shared state & normalised value or configuration & utility → caller \\\hline
svgPlot & 5119--5163 & Client / view & DOM state, selection, user command & HTML/SVG, status and download command & view state → render \\\hline
svgMessage & 5164--5164 & Client / view & DOM state, selection, user command & HTML/SVG, status and download command & view state → render \\\hline
extent & 5165--5165 & Shared model / utility & primitive values or shared state & normalised value or configuration & utility → caller \\\hline
legendFromMap & 5166--5166 & Shared model / utility & primitive values or shared state & normalised value or configuration & utility → caller \\\hline
catColour & 5167--5167 & Shared model / utility & primitive values or shared state & normalised value or configuration & utility → caller \\\hline
sopColour & 5168--5168 & Controller / service & RUNS, profile, history or chart inputs & derived rows, findings, limits or calls & model → next service/view \\\hline
sopLabel & 5169--5169 & Controller / service & RUNS, profile, history or chart inputs & derived rows, findings, limits or calls & model → next service/view \\\hline
gRunWells & 5172--5172 & Controller / service & RUNS, profile, history or chart inputs & derived rows, findings, limits or calls & model → next service/view \\\hline
gTargets & 5173--5173 & Shared model / utility & primitive values or shared state & normalised value or configuration & utility → caller \\\hline
gCutLines & 5174--5180 & Shared model / utility & primitive values or shared state & normalised value or configuration & utility → caller \\\hline
gRunDate & 5183--5183 & Controller / service & RUNS, profile, history or chart inputs & derived rows, findings, limits or calls & model → next service/view \\\hline
gByDate & 5184--5184 & Shared model / utility & primitive values or shared state & normalised value or configuration & utility → caller \\\hline
gRunTicks & 5185--5190 & Controller / service & RUNS, profile, history or chart inputs & derived rows, findings, limits or calls & model → next service/view \\\hline
GRAPHS & 5191--5542 & Client / view & DOM state, selection, user command & HTML/SVG, status and download command & view state → render \\\hline
GRAPH\_METRICS & 5543--5543 & Client / view & DOM state, selection, user command & HTML/SVG, status and download command & view state → render \\\hline
GRAPH\_STATE & 5544--5544 & Client / view & DOM state, selection, user command & HTML/SVG, status and download command & view state → render \\\hline
graphContext & 5545--5550 & Client / view & DOM state, selection, user command & HTML/SVG, status and download command & view state → render \\\hline
renderGraphs & 5551--5580 & Client / view & DOM state, selection, user command & HTML/SVG, status and download command & view state → render \\\hline
graphFileBase & 5581--5581 & Client / view & DOM state, selection, user command & HTML/SVG, status and download command & view state → render \\\hline
graphSvgText & 5582--5583 & Client / view & DOM state, selection, user command & HTML/SVG, status and download command & view state → render \\\hline
graphPng & 5584--5591 & Client / view & DOM state, selection, user command & HTML/SVG, status and download command & view state → render \\\hline
graphAllZip & 5592--5612 & Client / view & DOM state, selection, user command & HTML/SVG, status and download command & view state → render \\\hline
wireGraphs & 5613--5625 & Client / view & DOM state, selection, user command & HTML/SVG, status and download command & view state → render \\\hline
markDirty & 5629--5629 & Shared model / utility & primitive values or shared state & normalised value or configuration & utility → caller \\\hline
currentTab & 5630--5630 & Client / view & DOM state, selection, user command & HTML/SVG, status and download command & view state → render \\\hline
currentReviewSheet & 5631--5631 & Client / view & DOM state, selection, user command & HTML/SVG, status and download command & view state → render \\\hline
REVIEW\_RENDERERS & 5632--5634 & Client / view & DOM state, selection, user command & HTML/SVG, status and download command & view state → render \\\hline
renderTab & 5635--5656 & Client / view & DOM state, selection, user command & HTML/SVG, status and download command & view state → render \\\hline
renderReviewSheetLazy & 5657--5661 & Client / view & DOM state, selection, user command & HTML/SVG, status and download command & view state → render \\\hline
ccLzss & 5676--5690 & Controller / service & RUNS, profile, history or chart inputs & derived rows, findings, limits or calls & model → next service/view \\\hline
ccDecodeArz & 5691--5704 & Controller / service & RUNS, profile, history or chart inputs & derived rows, findings, limits or calls & model → next service/view \\\hline
ccIxoTemperatureLog & 5705--5713 & Controller / service & RUNS, profile, history or chart inputs & derived rows, findings, limits or calls & model → next service/view \\\hline
ccIxoAcquisitions & 5716--5723 & Client / view & DOM state, selection, user command & HTML/SVG, status and download command & view state → render \\\hline
ccThermalSteps & 5726--5758 & Controller / service & RUNS, profile, history or chart inputs & derived rows, findings, limits or calls & model → next service/view \\\hline
ccTransitionPairs & 5766--5777 & Controller / service & RUNS, profile, history or chart inputs & derived rows, findings, limits or calls & model → next service/view \\\hline
ccTransitions & 5778--5804 & Controller / service & RUNS, profile, history or chart inputs & derived rows, findings, limits or calls & model → next service/view \\\hline
ccThin & 5807--5814 & Controller / service & RUNS, profile, history or chart inputs & derived rows, findings, limits or calls & model → next service/view \\\hline
ccPeakRates & 5815--5827 & Controller / service & RUNS, profile, history or chart inputs & derived rows, findings, limits or calls & model → next service/view \\\hline
ccThermalSummary & 5828--5836 & Controller / service & RUNS, profile, history or chart inputs & derived rows, findings, limits or calls & model → next service/view \\\hline
ccSetpoints & 5837--5840 & Controller / service & RUNS, profile, history or chart inputs & derived rows, findings, limits or calls & model → next service/view \\\hline
ccQsLog & 5843--5883 & Controller / service & RUNS, profile, history or chart inputs & derived rows, findings, limits or calls & model → next service/view \\\hline
ccCalibSummary & 5885--5895 & Controller / service & RUNS, profile, history or chart inputs & derived rows, findings, limits or calls & model → next service/view \\\hline
ccQuantSaturation & 5897--5902 & Controller / service & RUNS, profile, history or chart inputs & derived rows, findings, limits or calls & model → next service/view \\\hline
CC\_TZ & 5903--5903 & Controller / service & RUNS, profile, history or chart inputs & derived rows, findings, limits or calls & model → next service/view \\\hline
ccParseJavaDate & 5904--5908 & Controller / service & RUNS, profile, history or chart inputs & derived rows, findings, limits or calls & model → next service/view \\\hline
ccQ & 5911--5911 & Controller / service & RUNS, profile, history or chart inputs & derived rows, findings, limits or calls & model → next service/view \\\hline
ccMedian & 5912--5912 & Controller / service & RUNS, profile, history or chart inputs & derived rows, findings, limits or calls & model → next service/view \\\hline
ccStatic & 5913--5979 & Controller / service & RUNS, profile, history or chart inputs & derived rows, findings, limits or calls & model → next service/view \\\hline
ccDynamic & 5981--6030 & Controller / service & RUNS, profile, history or chart inputs & derived rows, findings, limits or calls & model → next service/view \\\hline
ccExtractOnLoad & 6032--6051 & Controller / service & RUNS, profile, history or chart inputs & derived rows, findings, limits or calls & model → next service/view \\\hline
ccProtocolKey & 6055--6058 & Controller / service & RUNS, profile, history or chart inputs & derived rows, findings, limits or calls & model → next service/view \\\hline
CC\_GROUPS & 6071--6075 & Controller / service & RUNS, profile, history or chart inputs & derived rows, findings, limits or calls & model → next service/view \\\hline
CC\_CHARTS & 6076--6260 & Controller / service & RUNS, profile, history or chart inputs & derived rows, findings, limits or calls & model → next service/view \\\hline
CC\_BY\_ID & 6261--6261 & Controller / service & RUNS, profile, history or chart inputs & derived rows, findings, limits or calls & model → next service/view \\\hline
CC\_DEFAULT\_CFG & 6262--6263 & Controller / service & RUNS, profile, history or chart inputs & derived rows, findings, limits or calls & model → next service/view \\\hline
ccCfg & 6267--6268 & Controller / service & RUNS, profile, history or chart inputs & derived rows, findings, limits or calls & model → next service/view \\\hline
ccSetCfg & 6269--6270 & Controller / service & RUNS, profile, history or chart inputs & derived rows, findings, limits or calls & model → next service/view \\\hline
CC\_C4 & 6273--6273 & Controller / service & RUNS, profile, history or chart inputs & derived rows, findings, limits or calls & model → next service/view \\\hline
ccPhase & 6274--6274 & Controller / service & RUNS, profile, history or chart inputs & derived rows, findings, limits or calls & model → next service/view \\\hline
ccI & 6276--6288 & Controller / service & RUNS, profile, history or chart inputs & derived rows, findings, limits or calls & model → next service/view \\\hline
ccWestgard & 6290--6300 & Controller / service & RUNS, profile, history or chart inputs & derived rows, findings, limits or calls & model → next service/view \\\hline
ccEWMA & 6301--6306 & Controller / service & RUNS, profile, history or chart inputs & derived rows, findings, limits or calls & model → next service/view \\\hline
ccCUSUM & 6307--6322 & Controller / service & RUNS, profile, history or chart inputs & derived rows, findings, limits or calls & model → next service/view \\\hline
ccTrend & 6324--6339 & Controller / service & RUNS, profile, history or chart inputs & derived rows, findings, limits or calls & model → next service/view \\\hline
ccP & 6340--6345 & Controller / service & RUNS, profile, history or chart inputs & derived rows, findings, limits or calls & model → next service/view \\\hline
ccU & 6346--6349 & Controller / service & RUNS, profile, history or chart inputs & derived rows, findings, limits or calls & model → next service/view \\\hline
ccC & 6350--6353 & Controller / service & RUNS, profile, history or chart inputs & derived rows, findings, limits or calls & model → next service/view \\\hline
ccXbarS & 6354--6365 & Controller / service & RUNS, profile, history or chart inputs & derived rows, findings, limits or calls & model → next service/view \\\hline
ccFunnel & 6366--6372 & Controller / service & RUNS, profile, history or chart inputs & derived rows, findings, limits or calls & model → next service/view \\\hline
CC\_LIMIT\_TEXT & 6375--6433 & Controller / service & RUNS, profile, history or chart inputs & derived rows, findings, limits or calls & model → next service/view \\\hline
CONSOLE\_DICT & 6434--6438 & Shared model / utility & primitive values or shared state & normalised value or configuration & utility → caller \\\hline
ccAlarmDir & 6439--6445 & Controller / service & RUNS, profile, history or chart inputs & derived rows, findings, limits or calls & model → next service/view \\\hline
CC\_STATE & 6450--6450 & Controller / service & RUNS, profile, history or chart inputs & derived rows, findings, limits or calls & model → next service/view \\\hline
CC\_COL & 6452--6452 & Controller / service & RUNS, profile, history or chart inputs & derived rows, findings, limits or calls & model → next service/view \\\hline
CC\_SEG\_PALETTE & 6456--6456 & Controller / service & RUNS, profile, history or chart inputs & derived rows, findings, limits or calls & model → next service/view \\\hline
ccSignal & 6458--6458 & Controller / service & RUNS, profile, history or chart inputs & derived rows, findings, limits or calls & model → next service/view \\\hline
ccPlatform & 6459--6459 & Controller / service & RUNS, profile, history or chart inputs & derived rows, findings, limits or calls & model → next service/view \\\hline
ccInstrumentKey & 6460--6460 & Controller / service & RUNS, profile, history or chart inputs & derived rows, findings, limits or calls & model → next service/view \\\hline
ccOperator & 6461--6461 & Controller / service & RUNS, profile, history or chart inputs & derived rows, findings, limits or calls & model → next service/view \\\hline
ccRowsFromRuns & 6463--6477 & Controller / service & RUNS, profile, history or chart inputs & derived rows, findings, limits or calls & model → next service/view \\\hline
ccAllRows & 6478--6481 & Controller / service & RUNS, profile, history or chart inputs & derived rows, findings, limits or calls & model → next service/view \\\hline
ccInstruments & 6482--6482 & Controller / service & RUNS, profile, history or chart inputs & derived rows, findings, limits or calls & model → next service/view \\\hline
ccRowsFor & 6483--6488 & Controller / service & RUNS, profile, history or chart inputs & derived rows, findings, limits or calls & model → next service/view \\\hline
ccVariants & 6490--6492 & Controller / service & RUNS, profile, history or chart inputs & derived rows, findings, limits or calls & model → next service/view \\\hline
ccValue & 6493--6495 & Controller / service & RUNS, profile, history or chart inputs & derived rows, findings, limits or calls & model → next service/view \\\hline
ccXAxis & 6496--6502 & Controller / service & RUNS, profile, history or chart inputs & derived rows, findings, limits or calls & model → next service/view \\\hline
ccSpec & 6503--6508 & Controller / service & RUNS, profile, history or chart inputs & derived rows, findings, limits or calls & model → next service/view \\\hline
ccCompute & 6509--6559 & Controller / service & RUNS, profile, history or chart inputs & derived rows, findings, limits or calls & model → next service/view \\\hline
ccShift & 6563--6574 & Controller / service & RUNS, profile, history or chart inputs & derived rows, findings, limits or calls & model → next service/view \\\hline
ccStatus & 6575--6600 & Controller / service & RUNS, profile, history or chart inputs & derived rows, findings, limits or calls & model → next service/view \\\hline
ccSpark & 6602--6612 & Controller / service & RUNS, profile, history or chart inputs & derived rows, findings, limits or calls & model → next service/view \\\hline
ccVersionMarks & 6613--6617 & Controller / service & RUNS, profile, history or chart inputs & derived rows, findings, limits or calls & model → next service/view \\\hline
ccChartSvg & 6618--6677 & Client / view & DOM state, selection, user command & HTML/SVG, status and download command & view state → render \\\hline
ccSChartSvg & 6678--6684 & Client / view & DOM state, selection, user command & HTML/SVG, status and download command & view state → render \\\hline
renderPanel & 6686--6726 & Client / view & DOM state, selection, user command & HTML/SVG, status and download command & view state → render \\\hline
renderCcDetail & 6727--6796 & Client / view & DOM state, selection, user command & HTML/SVG, status and download command & view state → render \\\hline
ccPng & 6797--6802 & Controller / service & RUNS, profile, history or chart inputs & derived rows, findings, limits or calls & model → next service/view \\\hline
ccChartRows & 6803--6810 & Controller / service & RUNS, profile, history or chart inputs & derived rows, findings, limits or calls & model → next service/view \\\hline
ccSvgToPng & 6811--6816 & Client / view & DOM state, selection, user command & HTML/SVG, status and download command & view state → render \\\hline
ccZipAll & 6820--6868 & Controller / service & RUNS, profile, history or chart inputs & derived rows, findings, limits or calls & model → next service/view \\\hline
ccHistoryJSON & 6870--6874 & Controller / service & RUNS, profile, history or chart inputs & derived rows, findings, limits or calls & model → next service/view \\\hline
ccHistoryCSV & 6875--6881 & Controller / service & RUNS, profile, history or chart inputs & derived rows, findings, limits or calls & model → next service/view \\\hline
ccMigrateControlKeys & 6883--6895 & Controller / service & RUNS, profile, history or chart inputs & derived rows, findings, limits or calls & model → next service/view \\\hline
ccImport & 6896--6904 & Controller / service & RUNS, profile, history or chart inputs & derived rows, findings, limits or calls & model → next service/view \\\hline
ccRemember & 6905--6905 & Controller / service & RUNS, profile, history or chart inputs & derived rows, findings, limits or calls & model → next service/view \\\hline
ccRestore & 6906--6906 & Controller / service & RUNS, profile, history or chart inputs & derived rows, findings, limits or calls & model → next service/view \\\hline
ccEvents & 6908--6919 & Controller / service & RUNS, profile, history or chart inputs & derived rows, findings, limits or calls & model → next service/view \\\hline
wirePanel & 6921--6938 & Client / view & DOM state, selection, user command & HTML/SVG, status and download command & view state → render \\\hline
effWol & 6952--6970 & Shared model / utility & primitive values or shared state & normalised value or configuration & utility → caller \\\hline
effSliwinProxy & 6971--6984 & Shared model / utility & primitive values or shared state & normalised value or configuration & utility → caller \\\hline
efficiencyFromSlope & 6991--6995 & Shared model / utility & primitive values or shared state & normalised value or configuration & utility → caller \\\hline
\$ & 7000--7000 & Shared model / utility & primitive values or shared state & normalised value or configuration & utility → caller \\\hline
\$\$ & 7001--7001 & Shared model / utility & primitive values or shared state & normalised value or configuration & utility → caller \\\hline
APP\_LIMITS & 7006--7006 & Shared model / utility & primitive values or shared state & normalised value or configuration & utility → caller \\\hline
APP\_STATE & 7007--7007 & Controller / service & RUNS, profile, history or chart inputs & derived rows, findings, limits or calls & model → next service/view \\\hline
appError & 7008--7015 & Shared model / utility & primitive values or shared state & normalised value or configuration & utility → caller \\\hline
APP\_PHASE\_TRANSITIONS & 7016--7016 & Shared model / utility & primitive values or shared state & normalised value or configuration & utility → caller \\\hline
setPhase & 7017--7026 & Shared model / utility & primitive values or shared state & normalised value or configuration & utility → caller \\\hline
appDiagnostics & 7027--7027 & Shared model / utility & primitive values or shared state & normalised value or configuration & utility → caller \\\hline
validateRunBoundary & 7028--7037 & Controller / service & RUNS, profile, history or chart inputs & derived rows, findings, limits or calls & model → next service/view \\\hline
validateHistoryBoundary & 7038--7045 & Shared model / utility & primitive values or shared state & normalised value or configuration & utility → caller \\\hline
MG & 7055--7055 & Client / view & DOM state, selection, user command & HTML/SVG, status and download command & view state → render \\\hline
MG\_GLYPH & 7056--7056 & Client / view & DOM state, selection, user command & HTML/SVG, status and download command & view state → render \\\hline
MG\_WORD & 7057--7057 & Client / view & DOM state, selection, user command & HTML/SVG, status and download command & view state → render \\\hline
mgBuildQueue & 7060--7177 & Client / view & DOM state, selection, user command & HTML/SVG, status and download command & view state → render \\\hline
mgItem & 7180--7180 & Client / view & DOM state, selection, user command & HTML/SVG, status and download command & view state → render \\\hline
MG\_TYPE\_LABEL & 7181--7181 & Client / view & DOM state, selection, user command & HTML/SVG, status and download command & view state → render \\\hline
mgChartTypeOptions & 7182--7186 & Client / view & DOM state, selection, user command & HTML/SVG, status and download command & view state → render \\\hline
mgChartAxisOptions & 7187--7194 & Client / view & DOM state, selection, user command & HTML/SVG, status and download command & view state → render \\\hline
mgStatusPaint & 7195--7204 & Client / view & DOM state, selection, user command & HTML/SVG, status and download command & view state → render \\\hline
mgRender & 7206--7284 & Client / view & DOM state, selection, user command & HTML/SVG, status and download command & view state → render \\\hline
mgSheet & 7287--7301 & Client / view & DOM state, selection, user command & HTML/SVG, status and download command & view state → render \\\hline
mgSheetPick & 7302--7307 & Client / view & DOM state, selection, user command & HTML/SVG, status and download command & view state → render \\\hline
mgSheetClose & 7308--7308 & Client / view & DOM state, selection, user command & HTML/SVG, status and download command & view state → render \\\hline
mgItemSentence & 7309--7313 & Client / view & DOM state, selection, user command & HTML/SVG, status and download command & view state → render \\\hline
mgStage & 7321--7326 & Client / view & DOM state, selection, user command & HTML/SVG, status and download command & view state → render \\\hline
mgStageSet & 7327--7327 & Client / view & DOM state, selection, user command & HTML/SVG, status and download command & view state → render \\\hline
mgApplyChartView & 7328--7335 & Client / view & DOM state, selection, user command & HTML/SVG, status and download command & view state → render \\\hline
mgZoom & 7336--7342 & Client / view & DOM state, selection, user command & HTML/SVG, status and download command & view state → render \\\hline
mgPan & 7343--7343 & Client / view & DOM state, selection, user command & HTML/SVG, status and download command & view state → render \\\hline
MG\_COMMANDS & 7345--7376 & Client / view & DOM state, selection, user command & HTML/SVG, status and download command & view state → render \\\hline
mgCommand & 7377--7394 & Client / view & DOM state, selection, user command & HTML/SVG, status and download command & view state → render \\\hline
mgGo & 7395--7400 & Client / view & DOM state, selection, user command & HTML/SVG, status and download command & view state → render \\\hline
mgGoChart & 7401--7409 & Client / view & DOM state, selection, user command & HTML/SVG, status and download command & view state → render \\\hline
mgSetChartType & 7410--7413 & Client / view & DOM state, selection, user command & HTML/SVG, status and download command & view state → render \\\hline
mgSetChartAxis & 7414--7417 & Client / view & DOM state, selection, user command & HTML/SVG, status and download command & view state → render \\\hline
mgReadFiles & 7418--7424 & Client / view & DOM state, selection, user command & HTML/SVG, status and download command & view state → render \\\hline
mgChooseFiles & 7425--7429 & Client / view & DOM state, selection, user command & HTML/SVG, status and download command & view state → render \\\hline
mgLeaveTo & 7430--7430 & Client / view & DOM state, selection, user command & HTML/SVG, status and download command & view state → render \\\hline
MG\_PICT & 7443--7444 & Client / view & DOM state, selection, user command & HTML/SVG, status and download command & view state → render \\\hline
mgIt & 7445--7446 & Client / view & DOM state, selection, user command & HTML/SVG, status and download command & view state → render \\\hline
mgSafe & 7447--7447 & Client / view & DOM state, selection, user command & HTML/SVG, status and download command & view state → render \\\hline
mgTxt & 7448--7448 & Client / view & DOM state, selection, user command & HTML/SVG, status and download command & view state → render \\\hline
mgViewLoad & 7451--7506 & Client / view & DOM state, selection, user command & HTML/SVG, status and download command & view state → render \\\hline
mgViewSop & 7509--7575 & Client / view & DOM state, selection, user command & HTML/SVG, status and download command & view state → render \\\hline
mgViewResults & 7578--7656 & Client / view & DOM state, selection, user command & HTML/SVG, status and download command & view state → render \\\hline
mgPanelItems & 7659--7696 & Client / view & DOM state, selection, user command & HTML/SVG, status and download command & view state → render \\\hline
mgViewPanel & 7697--7719 & Client / view & DOM state, selection, user command & HTML/SVG, status and download command & view state → render \\\hline
mgGraphCtx & 7722--7726 & Client / view & DOM state, selection, user command & HTML/SVG, status and download command & view state → render \\\hline
mgViewGraphs & 7727--7755 & Client / view & DOM state, selection, user command & HTML/SVG, status and download command & view state → render \\\hline
mgViewReview & 7758--7788 & Client / view & DOM state, selection, user command & HTML/SVG, status and download command & view state → render \\\hline
mgViewIntegrity & 7791--7803 & Client / view & DOM state, selection, user command & HTML/SVG, status and download command & view state → render \\\hline
mgViewExport & 7806--7831 & Client / view & DOM state, selection, user command & HTML/SVG, status and download command & view state → render \\\hline
MG\_VIEWS & 7833--7842 & Client / view & DOM state, selection, user command & HTML/SVG, status and download command & view state → render \\\hline
mgScope & 7843--7846 & Client / view & DOM state, selection, user command & HTML/SVG, status and download command & view state → render \\\hline
mgScopeMark & 7847--7850 & Client / view & DOM state, selection, user command & HTML/SVG, status and download command & view state → render \\\hline
mgItemKey & 7851--7851 & Client / view & DOM state, selection, user command & HTML/SVG, status and download command & view state → render \\\hline
mgScopeFilter & 7854--7859 & Client / view & DOM state, selection, user command & HTML/SVG, status and download command & view state → render \\\hline
mgRefresh & 7860--7874 & Client / view & DOM state, selection, user command & HTML/SVG, status and download command & view state → render \\\hline
mgOpen & 7878--7882 & Client / view & DOM state, selection, user command & HTML/SVG, status and download command & view state → render \\\hline
mgClose & 7883--7885 & Client / view & DOM state, selection, user command & HTML/SVG, status and download command & view state → render \\\hline
mgToggle & 7886--7886 & Client / view & DOM state, selection, user command & HTML/SVG, status and download command & view state → render \\\hline
mgInit & 7887--7897 & Client / view & DOM state, selection, user command & HTML/SVG, status and download command & view state → render \\\hline
APP\_STATS & 7904--7904 & Shared model / utility & primitive values or shared state & normalised value or configuration & utility → caller \\\hline
appTaskStart & 7905--7909 & Shared model / utility & primitive values or shared state & normalised value or configuration & utility → caller \\\hline
appTaskStep & 7910--7910 & Shared model / utility & primitive values or shared state & normalised value or configuration & utility → caller \\\hline
appTaskEnd & 7911--7915 & Shared model / utility & primitive values or shared state & normalised value or configuration & utility → caller \\\hline
appTaskWrap & 7916--7919 & Shared model / utility & primitive values or shared state & normalised value or configuration & utility → caller \\\hline
APP\_TASK\_LABEL & 7920--7922 & Shared model / utility & primitive values or shared state & normalised value or configuration & utility → caller \\\hline
appHeapMB & 7923--7923 & Shared model / utility & primitive values or shared state & normalised value or configuration & utility → caller \\\hline
appStagedBytes & 7924--7924 & Repository / data access & File bytes, decoded rows or export request & stored model, checksum or file bytes & repository → model/export \\\hline
appNotice & 7925--7930 & Shared model / utility & primitive values or shared state & normalised value or configuration & utility → caller \\\hline
appCancel & 7931--7935 & Shared model / utility & primitive values or shared state & normalised value or configuration & utility → caller \\\hline
appReadSuperseded & 7936--7943 & Repository / data access & File bytes, decoded rows or export request & stored model, checksum or file bytes & repository → model/export \\\hline
APP\_PHASE\_TEXT & 7944--7945 & Shared model / utility & primitive values or shared state & normalised value or configuration & utility → caller \\\hline
appFmtMs & 7946--7946 & Shared model / utility & primitive values or shared state & normalised value or configuration & utility → caller \\\hline
appFmtMB & 7947--7947 & Shared model / utility & primitive values or shared state & normalised value or configuration & utility → caller \\\hline
appStatusSnapshot & 7948--7961 & Shared model / utility & primitive values or shared state & normalised value or configuration & utility → caller \\\hline
renderConsoleIdentity & 7962--7973 & Client / view & DOM state, selection, user command & HTML/SVG, status and download command & view state → render \\\hline
appStatusPaint & 7975--8022 & Shared model / utility & primitive values or shared state & normalised value or configuration & utility → caller \\\hline
appStatusReport & 8023--8028 & Shared model / utility & primitive values or shared state & normalised value or configuration & utility → caller \\\hline
appStatusInit & 8029--8040 & Shared model / utility & primitive values or shared state & normalised value or configuration & utility → caller \\\hline
RUNS & 8041--8041 & Controller / service & RUNS, profile, history or chart inputs & derived rows, findings, limits or calls & model → next service/view \\\hline
tag & 8045--8045 & Shared model / utility & primitive values or shared state & normalised value or configuration & utility → caller \\\hline
table & 8046--8054 & Client / view & DOM state, selection, user command & HTML/SVG, status and download command & view state → render \\\hline
download & 8055--8061 & Repository / data access & File bytes, decoded rows or export request & stored model, checksum or file bytes & repository → model/export \\\hline
websiteSource & 8062--8065 & Repository / data access & File bytes, decoded rows or export request & stored model, checksum or file bytes & repository → model/export \\\hline
makeWebsiteArchive & 8066--8072 & Shared model / utility & primitive values or shared state & normalised value or configuration & utility → caller \\\hline
downloadWebsiteArchive & 8073--8076 & Repository / data access & File bytes, decoded rows or export request & stored model, checksum or file bytes & repository → model/export \\\hline
csvq & 8077--8080 & Repository / data access & File bytes, decoded rows or export request & stored model, checksum or file bytes & repository → model/export \\\hline
toCSV & 8081--8084 & Repository / data access & File bytes, decoded rows or export request & stored model, checksum or file bytes & repository → model/export \\\hline
stamp & 8085--8085 & Shared model / utility & primitive values or shared state & normalised value or configuration & utility → caller \\\hline
baseName & 8086--8090 & Shared model / utility & primitive values or shared state & normalised value or configuration & utility → caller \\\hline
NAME\_ROLE\_RULES & 8103--8110 & Shared model / utility & primitive values or shared state & normalised value or configuration & utility → caller \\\hline
roleFromName & 8111--8115 & Shared model / utility & primitive values or shared state & normalised value or configuration & utility → caller \\\hline
assignRoles & 8116--8143 & Shared model / utility & primitive values or shared state & normalised value or configuration & utility → caller \\\hline
operatorOf & 8148--8148 & Shared model / utility & primitive values or shared state & normalised value or configuration & utility → caller \\\hline
sourceUidOf & 8149--8149 & Client / view & DOM state, selection, user command & HTML/SVG, status and download command & view state → render \\\hline
sourceCRCOf & 8150--8150 & Repository / data access & File bytes, decoded rows or export request & stored model, checksum or file bytes & repository → model/export \\\hline
sourceSHA256Of & 8151--8151 & Repository / data access & File bytes, decoded rows or export request & stored model, checksum or file bytes & repository → model/export \\\hline
runDurationMinutes & 8152--8155 & Controller / service & RUNS, profile, history or chart inputs & derived rows, findings, limits or calls & model → next service/view \\\hline
uniqueCurveCount & 8156--8158 & Controller / service & RUNS, profile, history or chart inputs & derived rows, findings, limits or calls & model → next service/view \\\hline
acquiredCurveCount & 8159--8159 & Client / view & DOM state, selection, user command & HTML/SVG, status and download command & view state → render \\\hline
wellRows & 8160--8190 & Shared model / utility & primitive values or shared state & normalised value or configuration & utility → caller \\\hline
allWellRows & 8191--8191 & Shared model / utility & primitive values or shared state & normalised value or configuration & utility → caller \\\hline
tmRows & 8192--8214 & Shared model / utility & primitive values or shared state & normalised value or configuration & utility → caller \\\hline
metaRows & 8215--8238 & Shared model / utility & primitive values or shared state & normalised value or configuration & utility → caller \\\hline
dyeChemistryOf & 8254--8263 & Shared model / utility & primitive values or shared state & normalised value or configuration & utility → caller \\\hline
rdmlSampleType & 8264--8269 & Repository / data access & File bytes, decoded rows or export request & stored model, checksum or file bytes & repository → model/export \\\hline
rdmlId & 8270--8270 & Repository / data access & File bytes, decoded rows or export request & stored model, checksum or file bytes & repository → model/export \\\hline
analysisGroups & 8271--8279 & Controller / service & RUNS, profile, history or chart inputs & derived rows, findings, limits or calls & model → next service/view \\\hline
curveAnalysisGroups & 8280--8280 & Controller / service & RUNS, profile, history or chart inputs & derived rows, findings, limits or calls & model → next service/view \\\hline
makeRDMLXML & 8282--8325 & Repository / data access & File bytes, decoded rows or export request & stored model, checksum or file bytes & repository → model/export \\\hline
makeRDML & 8326--8326 & Repository / data access & File bytes, decoded rows or export request & stored model, checksum or file bytes & repository → model/export \\\hline
makeQpcrWide & 8331--8343 & Shared model / utility & primitive values or shared state & normalised value or configuration & utility → caller \\\hline
makeQpcrWideTables & 8344--8351 & Client / view & DOM state, selection, user command & HTML/SVG, status and download command & view state → render \\\hline
makeBaselineCSV & 8355--8367 & Repository / data access & File bytes, decoded rows or export request & stored model, checksum or file bytes & repository → model/export \\\hline
makeRDES & 8368--8380 & Repository / data access & File bytes, decoded rows or export request & stored model, checksum or file bytes & repository → model/export \\\hline
makeRDESTables & 8381--8388 & Client / view & DOM state, selection, user command & HTML/SVG, status and download command & view state → render \\\hline
rScript & 8391--8429 & Shared model / utility & primitive values or shared state & normalised value or configuration & utility → caller \\\hline
pyScript & 8430--8468 & Shared model / utility & primitive values or shared state & normalised value or configuration & utility → caller \\\hline
BASE\_MACHINE\_METADATA & 8473--8473 & Repository / data access & File bytes, decoded rows or export request & stored model, checksum or file bytes & repository → model/export \\\hline
JSONLD\_STANDARDS & 8474--8480 & Shared model / utility & primitive values or shared state & normalised value or configuration & utility → caller \\\hline
metadataId & 8481--8484 & Repository / data access & File bytes, decoded rows or export request & stored model, checksum or file bytes & repository → model/export \\\hline
outputDistribution & 8485--8491 & Shared model / utility & primitive values or shared state & normalised value or configuration & utility → caller \\\hline
distributionsFor & 8492--8519 & Shared model / utility & primitive values or shared state & normalised value or configuration & utility → caller \\\hline
variablesFor & 8520--8527 & Shared model / utility & primitive values or shared state & normalised value or configuration & utility → caller \\\hline
datasetMetadata & 8528--8573 & Repository / data access & File bytes, decoded rows or export request & stored model, checksum or file bytes & repository → model/export \\\hline
pastedDatasetMetadata & 8574--8585 & Repository / data access & File bytes, decoded rows or export request & stored model, checksum or file bytes & repository → model/export \\\hline
machineMetadataDocument & 8586--8598 & Repository / data access & File bytes, decoded rows or export request & stored model, checksum or file bytes & repository → model/export \\\hline
machineMetadataJSON & 8599--8599 & Repository / data access & File bytes, decoded rows or export request & stored model, checksum or file bytes & repository → model/export \\\hline
renderMachineMetadata & 8600--8603 & Client / view & DOM state, selection, user command & HTML/SVG, status and download command & view state → render \\\hline
exportCatalogue & 8604--8640 & Repository / data access & File bytes, decoded rows or export request & stored model, checksum or file bytes & repository → model/export \\\hline
exportCatalogueCore & 8641--8693 & Repository / data access & File bytes, decoded rows or export request & stored model, checksum or file bytes & repository → model/export \\\hline
bundleZip & 8694--8737 & Repository / data access & File bytes, decoded rows or export request & stored model, checksum or file bytes & repository → model/export \\\hline
bundleReadme & 8738--8780 & Repository / data access & File bytes, decoded rows or export request & stored model, checksum or file bytes & repository → model/export \\\hline
renderLoaded & 8786--8816 & Client / view & DOM state, selection, user command & HTML/SVG, status and download command & view state → render \\\hline
renderDiagnostics & 8817--8833 & Client / view & DOM state, selection, user command & HTML/SVG, status and download command & view state → render \\\hline
renderExports & 8834--8873 & Client / view & DOM state, selection, user command & HTML/SVG, status and download command & view state → render \\\hline
REVIEW\_COLOURS & 8882--8888 & Controller / service & RUNS, profile, history or chart inputs & derived rows, findings, limits or calls & model → next service/view \\\hline
PLATE\_PASTELS & 8889--8889 & Shared model / utility & primitive values or shared state & normalised value or configuration & utility → caller \\\hline
CONTROL\_CATEGORY\_OPTIONS & 8890--8892 & Client / view & DOM state, selection, user command & HTML/SVG, status and download command & view state → render \\\hline
REVIEW\_STATE & 8893--8894 & Controller / service & RUNS, profile, history or chart inputs & derived rows, findings, limits or calls & model → next service/view \\\hline
runName & 8896--8896 & Controller / service & RUNS, profile, history or chart inputs & derived rows, findings, limits or calls & model → next service/view \\\hline
shortDate & 8897--8900 & Shared model / utility & primitive values or shared state & normalised value or configuration & utility → caller \\\hline
displayRunTick & 8901--8904 & Controller / service & RUNS, profile, history or chart inputs & derived rows, findings, limits or calls & model → next service/view \\\hline
runAt & 8905--8908 & Controller / service & RUNS, profile, history or chart inputs & derived rows, findings, limits or calls & model → next service/view \\\hline
setSelectItems & 8909--8916 & Client / view & DOM state, selection, user command & HTML/SVG, status and download command & view state → render \\\hline
analysisChoices & 8918--8928 & Controller / service & RUNS, profile, history or chart inputs & derived rows, findings, limits or calls & model → next service/view \\\hline
selectedWells & 8929--8933 & Shared model / utility & primitive values or shared state & normalised value or configuration & utility → caller \\\hline
physicalCurveRows & 8934--8950 & Controller / service & RUNS, profile, history or chart inputs & derived rows, findings, limits or calls & model → next service/view \\\hline
heatColour & 8951--8956 & Client / view & DOM state, selection, user command & HTML/SVG, status and download command & view state → render \\\hline
plateMapRows & 8957--8984 & Shared model / utility & primitive values or shared state & normalised value or configuration & utility → caller \\\hline
sampleComparisonRows & 8985--8998 & Shared model / utility & primitive values or shared state & normalised value or configuration & utility → caller \\\hline
reviewStatisticsRows & 8999--9018 & Controller / service & RUNS, profile, history or chart inputs & derived rows, findings, limits or calls & model → next service/view \\\hline
replicateReviewRows & 9019--9033 & Controller / service & RUNS, profile, history or chart inputs & derived rows, findings, limits or calls & model → next service/view \\\hline
BLANK\_ROLES & 9040--9040 & Shared model / utility & primitive values or shared state & normalised value or configuration & utility → caller \\\hline
orphanBlankRole & 9041--9047 & Controller / service & RUNS, profile, history or chart inputs & derived rows, findings, limits or calls & model → next service/view \\\hline
orphanOtherChannels & 9051--9063 & Controller / service & RUNS, profile, history or chart inputs & derived rows, findings, limits or calls & model → next service/view \\\hline
resultCq & 9072--9077 & Shared model / utility & primitive values or shared state & normalised value or configuration & utility → caller \\\hline
resultAnswers & 9078--9086 & Shared model / utility & primitive values or shared state & normalised value or configuration & utility → caller \\\hline
quantChannels & 9092--9099 & Shared model / utility & primitive values or shared state & normalised value or configuration & utility → caller \\\hline
orphanCurveCandidates & 9100--9192 & Controller / service & RUNS, profile, history or chart inputs & derived rows, findings, limits or calls & model → next service/view \\\hline
reviewOrphanWells & 9193--9201 & Controller / service & RUNS, profile, history or chart inputs & derived rows, findings, limits or calls & model → next service/view \\\hline
reviewForensicEventsCoreUncached & 9202--9260 & Controller / service & RUNS, profile, history or chart inputs & derived rows, findings, limits or calls & model → next service/view \\\hline
reviewRunRows & 9261--9279 & Controller / service & RUNS, profile, history or chart inputs & derived rows, findings, limits or calls & model → next service/view \\\hline
stripLotDate & 9280--9284 & Shared model / utility & primitive values or shared state & normalised value or configuration & utility → caller \\\hline
controlRootName & 9285--9288 & Controller / service & RUNS, profile, history or chart inputs & derived rows, findings, limits or calls & model → next service/view \\\hline
controlRootKey & 9289--9291 & Controller / service & RUNS, profile, history or chart inputs & derived rows, findings, limits or calls & model → next service/view \\\hline
controlProfileRoot & 9292--9302 & Controller / service & RUNS, profile, history or chart inputs & derived rows, findings, limits or calls & model → next service/view \\\hline
controlCategory & 9303--9324 & Controller / service & RUNS, profile, history or chart inputs & derived rows, findings, limits or calls & model → next service/view \\\hline
detectedControlRoots & 9325--9340 & Controller / service & RUNS, profile, history or chart inputs & derived rows, findings, limits or calls & model → next service/view \\\hline
renderControlCategoryEditor & 9341--9365 & Client / view & DOM state, selection, user command & HTML/SVG, status and download command & view state → render \\\hline
westgardReview & 9366--9383 & Controller / service & RUNS, profile, history or chart inputs & derived rows, findings, limits or calls & model → next service/view \\\hline
reviewControlSeries & 9384--9420 & Controller / service & RUNS, profile, history or chart inputs & derived rows, findings, limits or calls & model → next service/view \\\hline
instrumentRecordTables & 9421--9454 & Client / view & DOM state, selection, user command & HTML/SVG, status and download command & view state → render \\\hline
emptyReview & 9456--9458 & Controller / service & RUNS, profile, history or chart inputs & derived rows, findings, limits or calls & model → next service/view \\\hline
metricCard & 9459--9461 & Client / view & DOM state, selection, user command & HTML/SVG, status and download command & view state → render \\\hline
downloadCanvas & 9462--9466 & Repository / data access & File bytes, decoded rows or export request & stored model, checksum or file bytes & repository → model/export \\\hline
hashColour & 9467--9470 & Shared model / utility & primitive values or shared state & normalised value or configuration & utility → caller \\\hline
reviewBundleEntries & 9471--9492 & Controller / service & RUNS, profile, history or chart inputs & derived rows, findings, limits or calls & model → next service/view \\\hline
downloadReviewBundle & 9493--9496 & Controller / service & RUNS, profile, history or chart inputs & derived rows, findings, limits or calls & model → next service/view \\\hline
FINDING\_SEVERITY\_RANK & 9497--9497 & Shared model / utility & primitive values or shared state & normalised value or configuration & utility → caller \\\hline
findingTextKey & 9498--9498 & Shared model / utility & primitive values or shared state & normalised value or configuration & utility → caller \\\hline
findingEvidenceKey & 9510--9513 & Shared model / utility & primitive values or shared state & normalised value or configuration & utility → caller \\\hline
findingKey & 9514--9518 & Shared model / utility & primitive values or shared state & normalised value or configuration & utility → caller \\\hline
reviewDisplayFindings & 9519--9537 & Controller / service & RUNS, profile, history or chart inputs & derived rows, findings, limits or calls & model → next service/view \\\hline
reviewDisplayCounts & 9538--9546 & Controller / service & RUNS, profile, history or chart inputs & derived rows, findings, limits or calls & model → next service/view \\\hline
reviewEventsWithoutIntegrity & 9547--9549 & Controller / service & RUNS, profile, history or chart inputs & derived rows, findings, limits or calls & model → next service/view \\\hline
integrityMismatches & 9550--9564 & Repository / data access & File bytes, decoded rows or export request & stored model, checksum or file bytes & repository → model/export \\\hline
renderIntegrity & 9565--9587 & Client / view & DOM state, selection, user command & HTML/SVG, status and download command & view state → render \\\hline
renderReviewOverview & 9589--9618 & Client / view & DOM state, selection, user command & HTML/SVG, status and download command & view state → render \\\hline
plateColour & 9619--9625 & Client / view & DOM state, selection, user command & HTML/SVG, status and download command & view state → render \\\hline
renderReviewPlate & 9626--9662 & Client / view & DOM state, selection, user command & HTML/SVG, status and download command & view state → render \\\hline
drawPlateMapImage & 9663--9684 & Client / view & DOM state, selection, user command & HTML/SVG, status and download command & view state → render \\\hline
renderReviewComparison & 9685--9705 & Client / view & DOM state, selection, user command & HTML/SVG, status and download command & view state → render \\\hline
renderReviewStatistics & 9706--9742 & Client / view & DOM state, selection, user command & HTML/SVG, status and download command & view state → render \\\hline
renderReviewControls & 9743--9768 & Client / view & DOM state, selection, user command & HTML/SVG, status and download command & view state → render \\\hline
drawLeveyJennings & 9769--9802 & Client / view & DOM state, selection, user command & HTML/SVG, status and download command & view state → render \\\hline
renderReviewRuns & 9803--9841 & Client / view & DOM state, selection, user command & HTML/SVG, status and download command & view state → render \\\hline
curveSignal & 9842--9842 & Controller / service & RUNS, profile, history or chart inputs & derived rows, findings, limits or calls & model → next service/view \\\hline
curveRowsSelected & 9843--9862 & Controller / service & RUNS, profile, history or chart inputs & derived rows, findings, limits or calls & model → next service/view \\\hline
SERIES\_PALETTE & 9863--9863 & Shared model / utility & primitive values or shared state & normalised value or configuration & utility → caller \\\hline
CALL\_LINE & 9865--9865 & Shared model / utility & primitive values or shared state & normalised value or configuration & utility → caller \\\hline
curveKey & 9866--9866 & Controller / service & RUNS, profile, history or chart inputs & derived rows, findings, limits or calls & model → next service/view \\\hline
curveColour & 9867--9880 & Controller / service & RUNS, profile, history or chart inputs & derived rows, findings, limits or calls & model → next service/view \\\hline
renderReviewCurves & 9881--9910 & Client / view & DOM state, selection, user command & HTML/SVG, status and download command & view state → render \\\hline
drawAmplificationCurves & 9911--9939 & Client / view & DOM state, selection, user command & HTML/SVG, status and download command & view state → render \\\hline
bindCurveTip & 9940--9957 & Client / view & DOM state, selection, user command & HTML/SVG, status and download command & view state → render \\\hline
risingCurveRows & 9958--9960 & Controller / service & RUNS, profile, history or chart inputs & derived rows, findings, limits or calls & model → next service/view \\\hline
risingCurveDisplayRows & 9961--9972 & Controller / service & RUNS, profile, history or chart inputs & derived rows, findings, limits or calls & model → next service/view \\\hline
drawRisingCurveBundle & 9973--9993 & Client / view & DOM state, selection, user command & HTML/SVG, status and download command & view state → render \\\hline
renderReviewRising & 9994--10017 & Client / view & DOM state, selection, user command & HTML/SVG, status and download command & view state → render \\\hline
renderReviewInstrument & 10018--10055 & Client / view & DOM state, selection, user command & HTML/SVG, status and download command & view state → render \\\hline
instrumentZip & 10056--10062 & Repository / data access & File bytes, decoded rows or export request & stored model, checksum or file bytes & repository → model/export \\\hline
downloadCurveCSV & 10063--10063 & Controller / service & RUNS, profile, history or chart inputs & derived rows, findings, limits or calls & model → next service/view \\\hline
downloadCurveCSVCore & 10064--10071 & Controller / service & RUNS, profile, history or chart inputs & derived rows, findings, limits or calls & model → next service/view \\\hline
downloadForensicCSV & 10072--10079 & Controller / service & RUNS, profile, history or chart inputs & derived rows, findings, limits or calls & model → next service/view \\\hline
renderQualityReview & 10080--10083 & Client / view & DOM state, selection, user command & HTML/SVG, status and download command & view state → render \\\hline
renderQualityReviewLazy & 10084--10084 & Client / view & DOM state, selection, user command & HTML/SVG, status and download command & view state → render \\\hline
showReviewSheet & 10085--10090 & Client / view & DOM state, selection, user command & HTML/SVG, status and download command & view state → render \\\hline
wireQualityReview & 10091--10121 & Client / view & DOM state, selection, user command & HTML/SVG, status and download command & view state → render \\\hline
STAGED & 10126--10126 & Shared model / utility & primitive values or shared state & normalised value or configuration & utility → caller \\\hline
intake & 10127--10196 & Shared model / utility & primitive values or shared state & normalised value or configuration & utility → caller \\\hline
readStaged & 10197--10207 & Repository / data access & File bytes, decoded rows or export request & stored model, checksum or file bytes & repository → model/export \\\hline
readStagedCore & 10208--10249 & Repository / data access & File bytes, decoded rows or export request & stored model, checksum or file bytes & repository → model/export \\\hline
parseCqTable & 10257--10302 & Client / view & DOM state, selection, user command & HTML/SVG, status and download command & view state → render \\\hline
acceptPasted & 10315--10324 & Controller / service & RUNS, profile, history or chart inputs & derived rows, findings, limits or calls & model → next service/view \\\hline
refreshAll & 10329--10336 & Repository / data access & File bytes, decoded rows or export request & stored model, checksum or file bytes & repository → model/export \\\hline
showTab & 10337--10343 & Client / view & DOM state, selection, user command & HTML/SVG, status and download command & view state → render \\\hline
boot & 10344--10401 & Client / view & DOM state, selection, user command & HTML/SVG, status and download command & view state → render \\\hline
SAMPLE\_TYPE\_LABEL & 10411--10416 & Shared model / utility & primitive values or shared state & normalised value or configuration & utility → caller \\\hline
derivedSettings & 10417--10572 & Controller / service & RUNS, profile, history or chart inputs & derived rows, findings, limits or calls & model → next service/view \\\hline
derivedSettingsFor & 10574--10577 & Controller / service & RUNS, profile, history or chart inputs & derived rows, findings, limits or calls & model → next service/view \\\hline
renderDerived & 10578--10591 & Client / view & DOM state, selection, user command & HTML/SVG, status and download command & view state → render \\\hline
rdmlEntries & 10645--10648 & Repository / data access & File bytes, decoded rows or export request & stored model, checksum or file bytes & repository → model/export \\\hline
rdmlIsContainer & 10649--10651 & Repository / data access & File bytes, decoded rows or export request & stored model, checksum or file bytes & repository → model/export \\\hline
rdmlKids & 10656--10660 & Repository / data access & File bytes, decoded rows or export request & stored model, checksum or file bytes & repository → model/export \\\hline
rdmlKid & 10661--10661 & Repository / data access & File bytes, decoded rows or export request & stored model, checksum or file bytes & repository → model/export \\\hline
rdmlText & 10662--10662 & Repository / data access & File bytes, decoded rows or export request & stored model, checksum or file bytes & repository → model/export \\\hline
rdmlNum & 10663--10666 & Repository / data access & File bytes, decoded rows or export request & stored model, checksum or file bytes & repository → model/export \\\hline
rdmlAttrId & 10667--10667 & Repository / data access & File bytes, decoded rows or export request & stored model, checksum or file bytes & repository → model/export \\\hline
RDML\_SAMPLE\_ROLE & 10670--10670 & Repository / data access & File bytes, decoded rows or export request & stored model, checksum or file bytes & repository → model/export \\\hline
rdmlParse & 10671--10832 & Controller / service & RUNS, profile, history or chart inputs & derived rows, findings, limits or calls & model → next service/view \\\hline
runProcessingState & 10851--10868 & Controller / service & RUNS, profile, history or chart inputs & derived rows, findings, limits or calls & model → next service/view \\\hline
runIsInstrumentAnalysed & 10869--10871 & Controller / service & RUNS, profile, history or chart inputs & derived rows, findings, limits or calls & model → next service/view \\\hline
sopProcessingApplicability & 10876--10889 & Controller / service & RUNS, profile, history or chart inputs & derived rows, findings, limits or calls & model → next service/view \\\hline
rdmlStem & 10896--10898 & Repository / data access & File bytes, decoded rows or export request & stored model, checksum or file bytes & repository → model/export \\\hline
rdmlKeyOf & 10899--10901 & Repository / data access & File bytes, decoded rows or export request & stored model, checksum or file bytes & repository → model/export \\\hline
applyRawFilePriority & 10902--10919 & Shared model / utility & primitive values or shared state & normalised value or configuration & utility → caller \\\hline
runIsSuperseded & 10920--10920 & Controller / service & RUNS, profile, history or chart inputs & derived rows, findings, limits or calls & model → next service/view \\\hline
PDF\_PAGE & 10943--10943 & Shared model / utility & primitive values or shared state & normalised value or configuration & utility → caller \\\hline
PDF\_WINANSI & 10946--10949 & Shared model / utility & primitive values or shared state & normalised value or configuration & utility → caller \\\hline
PDF\_SPELL & 10950--10953 & Shared model / utility & primitive values or shared state & normalised value or configuration & utility → caller \\\hline
pdfStr & 10954--10969 & Shared model / utility & primitive values or shared state & normalised value or configuration & utility → caller \\\hline
pdfWidth & 10970--10987 & Shared model / utility & primitive values or shared state & normalised value or configuration & utility → caller \\\hline
pdfFit & 10988--10993 & Shared model / utility & primitive values or shared state & normalised value or configuration & utility → caller \\\hline
pdfWrap & 10994--11003 & Shared model / utility & primitive values or shared state & normalised value or configuration & utility → caller \\\hline
MISSY\_TEXT & 11007--11007 & Shared model / utility & primitive values or shared state & normalised value or configuration & utility → caller \\\hline
MISSY\_COLORS & 11008--11008 & Shared model / utility & primitive values or shared state & normalised value or configuration & utility → caller \\\hline
missyEmblem & 11010--11016 & Shared model / utility & primitive values or shared state & normalised value or configuration & utility → caller \\\hline
missyWords & 11017--11020 & Shared model / utility & primitive values or shared state & normalised value or configuration & utility → caller \\\hline
missySvg & 11021--11024 & Client / view & DOM state, selection, user command & HTML/SVG, status and download command & view state → render \\\hline
missyMarkWithSvg & 11025--11031 & Client / view & DOM state, selection, user command & HTML/SVG, status and download command & view state → render \\\hline
pdfDoc & 11032--11101 & Shared model / utility & primitive values or shared state & normalised value or configuration & utility → caller \\\hline
svgToJpeg & 11104--11120 & Client / view & DOM state, selection, user command & HTML/SVG, status and download command & view state → render \\\hline
SEL\_STATE & 11126--11126 & Controller / service & RUNS, profile, history or chart inputs & derived rows, findings, limits or calls & model → next service/view \\\hline
selRunIndex & 11127--11130 & Controller / service & RUNS, profile, history or chart inputs & derived rows, findings, limits or calls & model → next service/view \\\hline
selGraphContext & 11131--11136 & Client / view & DOM state, selection, user command & HTML/SVG, status and download command & view state → render \\\hline
selImageItems & 11137--11144 & Shared model / utility & primitive values or shared state & normalised value or configuration & utility → caller \\\hline
selDataItems & 11145--11183 & Shared model / utility & primitive values or shared state & normalised value or configuration & utility → caller \\\hline
SEL\_OPTION\_RULES & 11184--11190 & Client / view & DOM state, selection, user command & HTML/SVG, status and download command & view state → render \\\hline
selChartItems & 11191--11205 & Shared model / utility & primitive values or shared state & normalised value or configuration & utility → caller \\\hline
selCatalogue & 11206--11206 & Shared model / utility & primitive values or shared state & normalised value or configuration & utility → caller \\\hline
selItem & 11207--11207 & Shared model / utility & primitive values or shared state & normalised value or configuration & utility → caller \\\hline
selMatch & 11208--11218 & Shared model / utility & primitive values or shared state & normalised value or configuration & utility → caller \\\hline
selSelected & 11219--11222 & Shared model / utility & primitive values or shared state & normalised value or configuration & utility → caller \\\hline
selMatchAll & 11224--11228 & Shared model / utility & primitive values or shared state & normalised value or configuration & utility → caller \\\hline
selAdd & 11229--11229 & Shared model / utility & primitive values or shared state & normalised value or configuration & utility → caller \\\hline
selRemove & 11230--11230 & Shared model / utility & primitive values or shared state & normalised value or configuration & utility → caller \\\hline
selClear & 11231--11231 & Shared model / utility & primitive values or shared state & normalised value or configuration & utility → caller \\\hline
sopLab & 11237--11240 & Controller / service & RUNS, profile, history or chart inputs & derived rows, findings, limits or calls & model → next service/view \\\hline
sopLabName & 11241--11241 & Controller / service & RUNS, profile, history or chart inputs & derived rows, findings, limits or calls & model → next service/view \\\hline
sopAnalystName & 11242--11242 & Controller / service & RUNS, profile, history or chart inputs & derived rows, findings, limits or calls & model → next service/view \\\hline
selCsvTemplate & 11245--11250 & Repository / data access & File bytes, decoded rows or export request & stored model, checksum or file bytes & repository → model/export \\\hline
selApplyCsv & 11251--11273 & Repository / data access & File bytes, decoded rows or export request & stored model, checksum or file bytes & repository → model/export \\\hline
selZipEntries & 11276--11295 & Repository / data access & File bytes, decoded rows or export request & stored model, checksum or file bytes & repository → model/export \\\hline
selReadme & 11296--11319 & Repository / data access & File bytes, decoded rows or export request & stored model, checksum or file bytes & repository → model/export \\\hline
selWithExportNames & 11323--11327 & Repository / data access & File bytes, decoded rows or export request & stored model, checksum or file bytes & repository → model/export \\\hline
selDownloadZip & 11328--11332 & Repository / data access & File bytes, decoded rows or export request & stored model, checksum or file bytes & repository → model/export \\\hline
selReportTitle & 11340--11342 & Shared model / utility & primitive values or shared state & normalised value or configuration & utility → caller \\\hline
selTableBlock & 11343--11388 & Client / view & DOM state, selection, user command & HTML/SVG, status and download command & view state → render \\\hline
selBuildReport & 11389--11451 & Client / view & DOM state, selection, user command & HTML/SVG, status and download command & view state → render \\\hline
selDownloadReport & 11452--11456 & Repository / data access & File bytes, decoded rows or export request & stored model, checksum or file bytes & repository → model/export \\\hline
CMD & 11463--11463 & Shared model / utility & primitive values or shared state & normalised value or configuration & utility → caller \\\hline
CMD\_HELP & 11464--11488 & Shared model / utility & primitive values or shared state & normalised value or configuration & utility → caller \\\hline
cmdEl & 11489--11489 & Shared model / utility & primitive values or shared state & normalised value or configuration & utility → caller \\\hline
cmdPrint & 11490--11496 & Shared model / utility & primitive values or shared state & normalised value or configuration & utility → caller \\\hline
cmdTable & 11497--11501 & Client / view & DOM state, selection, user command & HTML/SVG, status and download command & view state → render \\\hline
cmdStatusLine & 11502--11504 & Shared model / utility & primitive values or shared state & normalised value or configuration & utility → caller \\\hline
selOptionRule & 11505--11508 & Client / view & DOM state, selection, user command & HTML/SVG, status and download command & view state → render \\\hline
selValidateOption & 11509--11515 & Client / view & DOM state, selection, user command & HTML/SVG, status and download command & view state → render \\\hline
selParseCommand & 11516--11523 & Controller / service & RUNS, profile, history or chart inputs & derived rows, findings, limits or calls & model → next service/view \\\hline
selApplyOptions & 11524--11524 & Client / view & DOM state, selection, user command & HTML/SVG, status and download command & view state → render \\\hline
cmdExec & 11525--11582 & Shared model / utility & primitive values or shared state & normalised value or configuration & utility → caller \\\hline
cmdRunLine & 11583--11590 & Controller / service & RUNS, profile, history or chart inputs & derived rows, findings, limits or calls & model → next service/view \\\hline
cmdRunScript & 11591--11599 & Controller / service & RUNS, profile, history or chart inputs & derived rows, findings, limits or calls & model → next service/view \\\hline
cmdScriptToggle & 11600--11604 & Shared model / utility & primitive values or shared state & normalised value or configuration & utility → caller \\\hline
cmdOpen & 11605--11614 & Shared model / utility & primitive values or shared state & normalised value or configuration & utility → caller \\\hline
cmdClose & 11615--11615 & Shared model / utility & primitive values or shared state & normalised value or configuration & utility → caller \\\hline
cmdToggle & 11616--11616 & Shared model / utility & primitive values or shared state & normalised value or configuration & utility → caller \\\hline
cmdInit & 11617--11639 & Shared model / utility & primitive values or shared state & normalised value or configuration & utility → caller \\\hline
AST & 11646--11646 & Shared model / utility & primitive values or shared state & normalised value or configuration & utility → caller \\\hline
astEl & 11647--11647 & Shared model / utility & primitive values or shared state & normalised value or configuration & utility → caller \\\hline
astRender & 11648--11680 & Client / view & DOM state, selection, user command & HTML/SVG, status and download command & view state → render \\\hline
astCount & 11681--11685 & Shared model / utility & primitive values or shared state & normalised value or configuration & utility → caller \\\hline
astOutput & 11686--11689 & Shared model / utility & primitive values or shared state & normalised value or configuration & utility → caller \\\hline
astBuild & 11690--11700 & Client / view & DOM state, selection, user command & HTML/SVG, status and download command & view state → render \\\hline
astOpen & 11701--11701 & Shared model / utility & primitive values or shared state & normalised value or configuration & utility → caller \\\hline
astClose & 11702--11702 & Shared model / utility & primitive values or shared state & normalised value or configuration & utility → caller \\\hline
astToggle & 11703--11703 & Shared model / utility & primitive values or shared state & normalised value or configuration & utility → caller \\\hline
astInit & 11704--11728 & Shared model / utility & primitive values or shared state & normalised value or configuration & utility → caller \\\hline
modesInit & 11729--11733 & Client / view & DOM state, selection, user command & HTML/SVG, status and download command & view state → render \\\hline
\bottomrule\end{longtable}\end{landscape}



\section{A3 architecture plates}
The following landscape plates provide larger versions of the composition, intake, review and export views for presentations. They are generated separately so the A3 geometry is preserved.
The A3 landscape design plates are delivered separately at \texttt{software\_design\_views\_A3\_landscape.pdf}; the preceding figures are the embedded schematic views.

